\chapter{الطريقة الدائرية ومدخل إلى غولدباخ ووارينغ}
\label{ch:circle-method-goldbach-waring}

\section{من مسألة جمعيّة إلى معامل فورييه}

تظهر كثير من المسائل الجمعية في نظرية الأعداد في صورة سؤال بسيط ظاهريًا: كم طريقة يمكن بها كتابة عدد صحيح كبير على صورة مجموع عناصر من مجموعة حسابية معينة؟ قد تكون العناصر قوى صحيحة، أو أعدادًا أولية، أو قيمًا لدالة كثيرة حدود. الصعوبة ليست في كتابة معادلة التمثيل، بل في استخراج عدد حلولها مع حد رئيس مضبوط وخطأ أصغر منه.

الفكرة المركزية للطريقة الدائرية هي تحويل شرط المساواة الصحيح إلى شرط تعامد فورييري. نستخدم طوال الفصل التطبيع
\[
e(t)=e^{2\pi i t}.
\]

\begin{proposition}[تعامد الدوال الأسية]
\label{prop:circle-orthogonality}
لكل عدد صحيح $m$ لدينا
\[
\int_0^1 e(\alpha m)\,d\alpha
=
\begin{cases}
1,&m=0,\\
0,&m\ne0.
\end{cases}
\]
\end{proposition}

\begin{proof}
إذا كان $m=0$ فالمكامل يساوي $1$. وإذا كان $m\ne0$ فإن
\[
\int_0^1 e(\alpha m)\,d\alpha
=
\left[\frac{e(\alpha m)}{2\pi i m}\right]_{0}^{1}
=
\frac{e(m)-1}{2\pi i m}
=0,
\]
لأن $m$ عدد صحيح ومن ثم $e(m)=1$.
\end{proof}

هذه الهوية هي البوابة التي تدخل منها مسائل التمثيل إلى التحليل. لتكن $\mathcal A$ مجموعة من الأعداد الصحيحة، ولتكن $w(n)$ أوزانًا ذات دعم منتهٍ، ونعرّف
\[
F(\alpha)=\sum_n w(n)e(\alpha n).
\]
عند رفع $F$ إلى القوة $s$ نحصل على
\[
F(\alpha)^s
=
\sum_{n_1,\ldots,n_s}
 w(n_1)\cdots w(n_s)
 e\!\left(\alpha(n_1+\cdots+n_s)\right).
\]

\begin{theorem}[هوية عد التمثيلات]
\label{thm:circle-representation-identity}
إذا عرّفنا عدد التمثيلات الموزون بـ
\[
R_s(N)
=
\sum_{n_1+\cdots+n_s=N}
 w(n_1)\cdots w(n_s),
\]
فإن
\[
R_s(N)
=
\int_0^1F(\alpha)^s e(-N\alpha)\,d\alpha.
\]
\end{theorem}

\begin{proof}
بالتوسيع والتبديل بين الجمع المنتهي والتكامل نجد
\[
\int_0^1F(\alpha)^s e(-N\alpha)\,d\alpha
=
\sum_{n_1,\ldots,n_s}
 w(n_1)\cdots w(n_s)
 \int_0^1
 e\!\left(\alpha(n_1+\cdots+n_s-N)\right)d\alpha.
\]
يختار التعامد بالضبط الحدود التي تحقق
$n_1+\cdots+n_s=N$، فتنتج الصيغة المطلوبة.
\end{proof}

\begin{remark}
هذه الهوية نتيجة دقيقة وليست تقريبًا. كل التقريبات اللاحقة تدخل عند تحليل التكامل، لا عند تحويل مسألة التمثيل إلى تكامل.
\end{remark}

\section{تشريح الدائرة: الأقواس الكبرى والأقواس الصغرى}

لا يكون المجموع الأسي $F(\alpha)$ بالحجم نفسه على الدائرة كلها. عندما تكون $\alpha$ قريبة من كسر ذي مقام صغير، تتراصف الأطوار جزئيًا وتظهر بنية حسابية محلية قوية. أما بعيدًا عن هذه الكسور، فيُنتظر أن يؤدي التذبذب إلى إلغاء كبير.

في نموذج وارينغ نضع
\[
P=N^{1/k},
\qquad
Q=P^{\eta},
\qquad
0<\eta<\frac{1}{4k},
\]
حيث $\eta$ ثابت صغير. لكل كسر مختزل $a/q$ مع $1\le q\le Q$ نعرّف القوس
\[
\mathfrak M(q,a)
=
\left\{\alpha\in[0,1):
\left|\alpha-\frac aq\right|
\le \frac{Q}{qN}
\right\}.
\]
ثم نضع
\[
\mathfrak M
=
\bigcup_{1\le q\le Q}
\bigcup_{\substack{1\le a\le q\\(a,q)=1}}
\mathfrak M(q,a),
\qquad
\mathfrak m=[0,1)\setminus\mathfrak M.
\]
نسمي $\mathfrak M$ مجموعة الأقواس الكبرى و$\mathfrak m$ مجموعة الأقواس الصغرى.

وبذلك تنقسم هوية التمثيل إلى
\[
R_s(N)
=
\int_{\mathfrak M}F(\alpha)^s e(-N\alpha)\,d\alpha
+
\int_{\mathfrak m}F(\alpha)^s e(-N\alpha)\,d\alpha.
\]

هذا الفصل ليس تقسيمًا شكليًا فقط. القسم الأول يحمل الحد الرئيس المتوقع، بينما يجب إثبات أن القسم الثاني أصغر منه. لذا تقوم الطريقة الدائرية على مهمتين مختلفتين:

\begin{enumerate}
\item بناء تقريب دقيق للمجموع الأسي قرب الكسور ذات المقامات الصغيرة.
\item إثبات إلغاء كافٍ على بقية الدائرة.
\end{enumerate}

\section{البنية المحلية والأرخميدية للحد الرئيس}

إذا كتبنا على قوس كبير
\[
\alpha=\frac aq+\beta,
\]
فإن المجموع الأسي يتفكك نموذجيًا إلى جزء منتهٍ يلتقط الحساب بترديد $q$، وجزء مستمر يلتقط السلوك الحجمي. في أبسط صورة إرشادية يظهر تقريب من النمط
\[
F\!\left(\frac aq+\beta\right)
\approx
q^{-1}S(q,a)V(\beta).
\]
هنا $S(q,a)$ مجموع محلي منتهٍ، و$V(\beta)$ تقريب أرخميدي مستمر.

بعد رفع هذا التقريب إلى القوة $s$، ثم الجمع على $a$ و$q$ والتكامل في $\beta$، تنفصل المساهمة إلى عاملين:
\[
\mathfrak S(N)\mathfrak J(N).
\]
يسمى $\mathfrak S(N)$ السلسلة المفردة، وهي تجمع الكثافات المحلية عند جميع الأوليات، ويسمى $\mathfrak J(N)$ التكامل المفرد، وهو العامل الأرخميدي الذي يقيس حجم منطقة الحلول الحقيقية.

هذا الفصل بين المحلي والأرخميدي من أهم مكاسب الطريقة الدائرية. فهو يوضح أن وجود حلول كثيرة يتطلب شرطين مختلفين:

\begin{enumerate}
\item عدم وجود عائق توافق محلي عند أي أولي.
\item وجود كتلة حقيقية موجبة للحلول بعد القياس المناسب.
\end{enumerate}

\begin{remark}[حدود الادعاء]
ظهور الرمزين $\mathfrak S$ و$\mathfrak J$ لا يثبت وحده صيغة تقاربية. لا بد من ضبط خطأ التقريب على الأقواس الكبرى، وإثبات تقارب السلسلة المفردة، وإيجابيتها عند الحاجة، ثم السيطرة على الأقواس الصغرى.
\end{remark}

\section{مدخل وارينغ}

لعدد صحيح ثابت $k\ge2$ نعرّف مجموع فايل
\[
f_k(\alpha;P)
=
\sum_{1\le x\le P}e(\alpha x^k),
\qquad P=N^{1/k}.
\]
وإذا كان $r_{s,k}(N)$ عدد حلول
\[
N=x_1^k+\cdots+x_s^k,
\qquad x_j\in\mathbb N,
\]
فإن مبرهنة~\ref{thm:circle-representation-identity} تعطي
\[
r_{s,k}(N)
=
\int_0^1 f_k(\alpha;P)^s e(-N\alpha)\,d\alpha.
\]

على الأقواس الكبرى تقود المجاميع الكاملة للقوى والتقريب المستمر إلى السلسلة والتكامل المفردين
\[
\mathfrak S_{s,k}(N),
\qquad
\mathfrak J_{s,k}(N).
\]
أما على الأقواس الصغرى فيحتاج البرهان إلى ادخار من رتبة قوة. المدخل المقتبس المعتمد في هذا الفصل هو وجود $\delta=\delta(k,s)>0$، عندما يكون $s$ كبيرًا بما يكفي بدلالة $k$، بحيث
\[
\int_{\mathfrak m}|f_k(\alpha;P)|^s\,d\alpha
\ll P^{s-k-\delta}.
\]
هذه نتيجة مقتبسة مركبة، وليست مبرهنة مثبتة هنا. تطوير أدوات المجاميع الأسية العامة، ومنها طريقة فان دير كوربوت، مؤجل إلى الفصل التالي.

وتقود البنية الكلاسيكية، تحت فروض العدد الكافي من المتغيرات وعدم وجود عائق محلي، إلى صيغة من الشكل
\[
r_{s,k}(N)
=
\mathfrak S_{s,k}(N)\mathfrak J_{s,k}(N)
+o\!\left(N^{s/k-1}\right).
\]
نستخدم هذه الصيغة بوصفها نتيجة كلاسيكية مقتبسة ومشروحة، ولا ندعي في هذا الفصل أفضل حد حديث لعدد المتغيرات.

\section{مدخل غولدباخ الثلاثي}

لدراسة تمثيل عدد فردي كمجموع ثلاثة أوليات نستخدم الوزن اللوغاريتمي
\[
S(\alpha)
=
\sum_{n\le N}\Lambda(n)e(\alpha n).
\]
ويكون عدد التمثيلات الموزون
\[
R_3(N)
=
\int_0^1S(\alpha)^3e(-N\alpha)\,d\alpha.
\]

قرب الكسور ذات المقامات الصغيرة يرتبط تحليل $S(\alpha)$ بتوزيع الأوليات في المتتاليات الحسابية. وعلى الأقواس الصغرى يلزم إلغاء عميق في المجاميع الأسية على الأوليات. لذلك لا يعاد في هذا الفصل برهان مبرهنة فينوغرادوف، بل تسجل بوصفها نتيجة مقتبسة:

\begin{theorem}[فينوغرادوف، مقتبس]
كل عدد صحيح فردي كبير بما يكفي هو مجموع ثلاثة أعداد أولية.
\end{theorem}

أكمل هلفغوت المسألة الثلاثية بإزالة قيد «كبير بما يكفي»:

\begin{theorem}[غولدباخ الثلاثي، هلفغوت، مقتبس]
كل عدد صحيح فردي أكبر من $5$ هو مجموع ثلاثة أعداد أولية.
\end{theorem}

البرهان الكامل لهذه النتيجة خارج نطاق الفصل؛ فهو يجمع تحليل الأقواس الكبرى والصغرى مع تقديرات صريحة وغربال كبير وتحقق حاسوبي منتهٍ. ويصنف الجزء الحاسوبي وحده `FINITE-VERIFIED`، لا برهانًا تحليليًا مستقلًا.

\section{غولدباخ الثنائية: الحد الفاصل بين المبرهنة والحدس}

تنص حدسية غولدباخ الثنائية على أن كل عدد زوجي أكبر من $2$ هو مجموع عددين أوليين. لا تثبت الطريقة الدائرية المعروفة هذه الحدسية حتى الآن. وهي تتنبأ بصيغة تقاربية ذات سلسلة مفردة محلية، لكن الأقواس الصغرى في المسألة الثنائية لا تترك هامشًا كافيًا بالتقنيات الكلاسيكية المتاحة.

لذلك يعتمد هذا الفصل التصنيف الصريح:

\[
\text{غولدباخ الثنائية} = \texttt{HYPOTHESIS / OPEN}.
\]

ولا يجوز تحويل التنبؤ الهاردي--ليتلوودي إلى مبرهنة، ولا الاستدلال من تحقق عددي منتهٍ إلى صحة عامة.

\section{ما الذي أنجزته الطريقة الدائرية؟}

الدرس البنيوي للطريقة الدائرية أوسع من أي تطبيق منفرد. فهي تبني جسرًا بين أربع طبقات:

\begin{enumerate}
\item معادلة جمعية صحيحة.
\item معامل فورييه دقيق لدالة توليد أسية.
\item كثافات محلية مجمعة في سلسلة مفردة.
\item حجم أرخميدي مجسد في تكامل مفرد.
\end{enumerate}

لكن هذا الجسر لا يكتمل إلا بإثبات الإلغاء على الأقواس الصغرى. ومن هنا تأتي ضرورة الفصل التالي: دراسة المجاميع الأسية وطريقة فان دير كوربوت بوصفهما الآلة التي تقيس التذبذب وتحوله إلى ادخار كمي.

\section*{حالة النتائج في هذا الفصل}

\begin{itemize}
\item تعامد الدوال الأسية: `IDENTITY / PROVED-HERE`.
\item هوية تكامل عدد التمثيلات: `IDENTITY / PROVED-HERE`.
\item تعريف الأقواس الكبرى والصغرى: `DEFINITION / PROVED-HERE`.
\item بنية السلسلة والتكامل المفردين: `PROVED-HERE / CITED INPUTS`.
\item الصيغة التقاربية الكلاسيكية في وارينغ: `CITED / EXPLAINED`.
\item مبرهنة فينوغرادوف: `CITED`.
\item مبرهنة هلفغوت: `CITED`.
\item التحقق الحاسوبي المنتهي ضمن الإكمال الحديث: `FINITE-VERIFIED`.
\item غولدباخ الثنائية: `HYPOTHESIS / OPEN`.
\end{itemize}
